\section{Ôn tập chương 5}
\Opensolutionfile{ans}[ans/ans-2-OTC5-DE1]
\subsection{Đề 1}
\caulc
% Câu 1%
\begin{ex}%[2H5H1-2]
	Cho mặt phẳng $(P)\colon x+2 y+3 z-1=0$. Véc-tơ nào dưới đây là một véc-tơ pháp tuyến của $(P)$?
	\choice
	{$\overrightarrow{n}_1=(1 ; 3 ;-1)$}
	{$\overrightarrow{n}_2=(2 ; 3 ;-1)$}
	{$\overrightarrow{n}_3=(1 ; 2 ;-1)$}
	{\True $\overrightarrow{n}_4=(1 ; 2 ; 3)$}
	\loigiai{
		$(P)$ có véc-tơ pháp tuyến là $\overrightarrow{n}_4=(1 ; 2 ; 3)$.
	}
\end{ex}
% Câu 2%
\begin{ex}%[2H5N1-3]
	Phương trình nào dưới đây là phương trình của mặt phẳng $(Oyz)$?
	\choice
	{$y=0$}
	{\True $x=0$}
	{$y-z=0$}
	{$z=0$}
	\loigiai{
		Phương trình của mặt phẳng $(O y z)$ là $x=0$.
	}
\end{ex}
% Câu 3%
\begin{ex}%[2H5H1-3]
	Phương trình nào dưới đây là phương trình mặt phẳng đi qua điểm $M(1 ; 2 ;-3)$ và có véc-tơ pháp tuyến $\overrightarrow{n}=(1 ;-2 ; 3)$?
	\choice
	{$x-2 y+3 z-12=0$}
	{$x-2 y-3 z+6=0$}
	{\True $x-2 y+3 z+12=0$}
	{$x-2 y-3 z-6=0$}
	\loigiai{
		Mặt phẳng đi qua điểm $M(1 ; 2 ;-3)$ và có véc-tơ pháp tuyến $\overrightarrow{n}=(1 ;-2 ; 3)$ nên phương trình mặt phẳng là $1(x-1)-2(y-2)+3(z+3)=0\Leftrightarrow x-2y+3z+12=0$.
	}
\end{ex}
% Câu 4%
\begin{ex}%[2H5H1-5]
	Cho mặt phẳng $(P)\colon 3 x+4 y+2 z+4=0$ và điểm $A(1 ;-2 ; 3)$. Khoảng cách từ $A$ đến $(P)$ bằng
	\choice
	{\True $\dfrac{5}{\sqrt{29}}$}
	{$\dfrac{5}{29}$}
	{$\dfrac{\sqrt{5}}{3}$}
	{$\dfrac{5}{9}$}
	\loigiai{
		Khoảng cách từ $A$ đến $(P)$ là 
		$$\mathrm{d}(A,(P))=\dfrac{\left| 3\cdot 1+ 4\cdot (-2)+2\cdot 3+4\right| }{\sqrt{3^2+4^2+2^2}}=\dfrac{5}{\sqrt{29}}.$$
	}
\end{ex}
% Câu 5%
\begin{ex}%[2H5H1-4]
	Cho ba mặt phẳng $(\alpha)\colon x+y+2 z+1=0$, $(\beta)\colon x+y-z+2=0$ và $(\gamma)\colon x-y+5=0$. Trong các mệnh đề sau, mệnh đề nào \textbf{sai}?
	\choice
	{$(\alpha) \perp(\beta)$}
	{$(\gamma) \perp(\beta)$}
	{\True $(\alpha) || (\beta)$}
	{$(\alpha) \perp(\gamma)$}
	\loigiai{
		Mặt phẳng $(\alpha)$ có véc-tơ	pháp tuyến $ \overrightarrow{n}_1=\left(1;1;2\right)$.\\
		Mặt phẳng $(\beta)$ có véc-tơ	pháp tuyến $ \overrightarrow{n}_2=\left(1;1;-1\right)$.\\
		Mặt phẳng $(\gamma)$ có véc-tơ	pháp tuyến $ \overrightarrow{n}_3=\left(1;-1;0\right)$.\\
		Khi đó ta có
		\begin{itemize}
			\item $\overrightarrow{n}_1\cdot \overrightarrow{n}_2=0 \Rightarrow (\alpha) \perp(\beta)$.
			\item $\overrightarrow{n}_2\cdot \overrightarrow{n}_3=0\Rightarrow (\gamma) \perp(\beta)$.
			\item $\overrightarrow{n}_1\cdot \overrightarrow{n}_3=0\Rightarrow (\alpha) \perp(\gamma)$.
		\end{itemize}
		
	}
\end{ex}
% Câu 6%
\begin{ex}%[2H5H2-2]
	Cho đường thẳng $d\colon \dfrac{x-2}{-1}=\dfrac{y-1}{2}=\dfrac{z+3}{1}$. Véc-tơ nào dưới đây là véc-tơ chỉ phương của đường thẳng $d$
	\choice
	{$\overrightarrow{u}_1=(2 ; 1 ;-3)$}
	{$\overrightarrow{u}_2=(-2 ;-1 ; 3)$}
	{\True $\overrightarrow{u}_3=(-1 ; 2 ; 1)$}
	{$\overrightarrow{u}_4=(-1 ; 2 ;-1)$}
	\loigiai{
		Đường thẳng $d$ có véc-tơ chỉ phương là $\overrightarrow{u}_3=(-1 ; 2 ; 1)$.
	}
\end{ex}
\begin{ex}%[2H5H2-4]
	Cho hai đường thẳng $d_1 \colon \heva{&x=1+2t \\&y=2+3t \\&z=3+4 t} \ (t \in \mathbb{R})$ và $d_2 \colon \heva{x=3+4 t' \\ y=5+6 t' \\ z=7+8 t'} \ (t' \in \mathbb{R})$.
	Trong các mệnh đề sau, mệnh đề nào đúng?
	\choice
	{$d_1$ và $d_2$ cắt nhau}
	{$d_1 \parallel d_2$}
	{\True $d_1 \equiv d_2$}
	{$d_1$ và $d_2$ chéo nhau}
	\loigiai{
		Đường thẳng $d_1$ có véc-tơ chỉ phương $\vec{u_1}=(2;3;4)$ và $M(1;2;3) \in d_1$.\\
		Đường thẳng $d_2$ có véc-tơ chỉ phương $\vec{u_2}=(4;6;8)$.\\
		Kiểm tra $M \in d_2$: thay $x=1$; $y=2$ và $z=3$ vào $d_2$ ta có hệ $$\heva{&1=3+4t'\\&2=5+6t'\\&3=7+8t'} \Leftrightarrow t'=-\dfrac{1}{2}.$$
		Suy ra $M \in d_2$. Mặt khác $\vec{u_1}$ cùng phương $\vec{u_2}$ nên $d_1$ trùng $d_2$.
	}
\end{ex}
\begin{ex}%[2H5H3-3]
	Phương trình mặt cầu tâm $I(2;1;-1)$, bán kính $R=2$ là
	\choice
	{$(x+2)^2+(y+1)^2+(z-1)^2=4$}
	{$(x+2)^2+(y+1)^2+(z-1)^2=2$}
	{$(x-2)^2+(y-1)^2+(z+1)^2=2$}
	{\True $(x-2)^2+(y-1)^2+(z+1)^2=4$}
	\loigiai{
		Phương trình mặt cầu tâm $I(2;1;-1)$, bán kính $R=2$ là	$(x-2)^2+(y-1)^2+(z+1)^2=4$.
	}
\end{ex}
\begin{ex}%[2H5H2-7]
	Góc giữa đường thẳng $d: \dfrac{x}{1}=\dfrac{y-1}{2}=\dfrac{z}{1}$ và mặt phẳng $(\alpha): 4x+3y+5z-4=0$ bằng
	\choice
	{$30^{\circ}$}
	{$120^{\circ}$}
	{\True $60^{\circ}$}
	{$150^{\circ}$}
	\loigiai{
		Đường thẳng $d$ có véc-tơ chỉ phương là $\vec{u}=(1;2;1)$.\\
		Mặt phẳng $(\alpha)$ có véc-tơ pháp tuyến là $\vec{n}=(4;3;5)$.\\
		$\sin(d,(\alpha))=\dfrac{|\vec{n}\cdot \vec{u}|}{|\vec{n}|\cdot|\vec{u}|}=\dfrac{|1\cdot 4+2\cdot 3+1\cdot 5|}{\sqrt{1^2+2^2+1^2}\sqrt{4^2+3^2+5^2}}=\dfrac{\sqrt{3}}{2} \Rightarrow (d,(\alpha))=60^{\circ}$.
	}
\end{ex}
\begin{ex}%[2H5V1-3]
	Trong không gian với hệ tọa độ $Oxyz$, mặt phẳng $(P)$ cắt ba trục $Ox$, $Oy$, $Oz$ tại $A$, $B$, $C$; trọng tâm của tam giác $ABC$ là $G(-1;-3;2)$. Phương trình mặt phẳng $(P)$ là
	\choice
	{$x+y-z-5=0$}
	{$2x-3y-z-1=0$}
	{$x+3y-2z+1=0$}
	{\True $6x+2y-3z+18=0$}
	\loigiai{
		$(P)$ cắt $Ox$, $Oy$, $Oz$ lần lượt tại $A(a;0;0)$; $B(0;b;0)$; $C(0;0;c)$; $(abc\neq 0)$.\\
		Tọa độ trọng tâm $G(-1;-3;2)$ suy ra: $\heva{&\dfrac{a}{3}=-1\\&\dfrac{b}{3}=-3\\&\dfrac{c}{3}=2}\Leftrightarrow\heva{&a=-3\\&b=-9\\&c=6.}$ \\
		Phương trình $(P)$ qua $A$, $B$, $C$ là $(P)\colon\dfrac{x}{-3}+\dfrac{y}{-9}+\dfrac{z}{6}=1\Leftrightarrow 6x+2y-3z+18=0$.}
\end{ex}
\begin{ex}%[2H5V2-3]
	Cho hai điểm $A(1;-2;-3)$, $B(-1;4;1)$ và đường thẳng $d \colon  \dfrac{x+2}{1}=\dfrac{y-2}{-1}=\dfrac{z+3}{2}$. Phương trình nào dưới đây là phương trình đường thẳng đi qua trung điểm của đoạn thẳng $AB$ và song song với $d$ ?
	\choice
	{$d' \colon  \dfrac{x}{1}=\dfrac{y-1}{1}=\dfrac{z+1}{2}$}
	{$d'\colon \dfrac{x}{1}=\dfrac{y-2}{-1}=\dfrac{z+2}{2}$}
	{\True $d'\colon \dfrac{x}{1}=\dfrac{y-1}{-1}=\dfrac{z+1}{2}$}
	{$d'\colon \dfrac{x-1}{1}=\dfrac{y-1}{-1}=\dfrac{z+1}{2}$}
	\loigiai{
		Ta có $I(0;1;-1)$ là trung điểm của $AB$ và $d$ có véc-tơ chỉ phương $\vec{u}_d=(1;-1;2)$.\\
		$d' \parallel d$ nên $d'$ nhận $\vec{u}_{d'}=\vec{u}_d=(1;-1;2)$ làm véc-tơ chỉ phương.\\
		Phương trình của $d'$ qua $I(0;1;-1)$ và nhận $\vec{u}_{d'}=(1;-1;2)$ làm véc-tơ chỉ phương là
		\[\dfrac{x}{1}=\dfrac{y-1}{-1}=\dfrac{z+1}{2}.\]
	}
\end{ex}


\begin{ex}%[2H5V3-3]
	Trong không gian với hệ tọa độ $Oxyz$, cho mặt cầu $(S)$ đi qua hai điểm $A(1;1;2)$, $B(3;0;1)$ và có tâm thuộc trục $Ox$. Phương trình mặt cầu $(S)$ là
	\choice
	{$(x+1)^{2}+y^{2}+z^{2}=5$}
	{$(x+1)^{2}+y^{2}+z^{2}=\sqrt{5}$}
	{\True $(x-1)^{2}+y^{2}+z^{2}=5$}
	{$(x-1)^{2}+y^{2}+z^{2}=\sqrt{5}$}
	\loigiai{
		$ I\in Ox $ nên gọi $ I(a;0;0) $ là tâm của mặt cầu. Vì mặt cầu qua hai điểm $A$, $B$ nên
		\[ IA=IB\Leftrightarrow\sqrt{(1-a)^2+1^2+2^2}=\sqrt{(3-a)^2+0^2+1^2} \Leftrightarrow a=1. \]
		Vậy mặt cầu có tâm $ I(1;0;0) $ và bán kính  $ R=IA=\sqrt{5}. $\\
		Phương trình mặt cầu $ (S) $ là  $ (x-1)^2+y^2+z^2=5 $.
	}
\end{ex}
\Closesolutionfile{ans}
\indapan{6}{ans/ans-2-OTC5-DE1}

\cauds
\Opensolutionfile{ans}[ans/ans-0-B15-DS]
\begin{ex}%[2H5V3-3]
	Phương trình tổng quát của mặt phẳng $(P)$ trong mỗi trường hợp sau là đúng hay sai?
	\choiceTF
		{\True $(P)$ đi qua điểm $M(-3;1;4)$ và có một véc-tơ pháp tuyến là $\vec{n}=(2;-4;1)$ là $2x-4y+z+6=0$}
		{$(P)$ đi qua điểm $N(2;-1;5)$ và có cặp véc-tơ chỉ phương là $\vec{u}_1=(1;-3;2)$ và $\vec{u}_2=(-3;4;1)$ là $11x+7y+5z+40=0$}
		{$(P)$ đi qua điểm $I(4;0;-7)$ và song song với mặt phẳng $(Q) \colon 2x-y-z-3=0$ là $2x+y-z-15=0$}
		{\True $(P)$ đi qua điểm $K(-4;9;2)$ và vuông góc với đường thẳng $\Delta \colon \dfrac{x-1}{2}=\dfrac{y}{1}=\dfrac{z-6}{5}$ là $2x+y+5z-11=0$}
	\loigiai{
	\begin{itemchoice}
			\itemch $(P)$ đi qua điểm $M(-3;1;4)$ và có một véc-tơ pháp tuyến là $\vec{n}=(2;-4;1)$ có phương trình là 
			$$2\cdot (x+3)-4\cdot(y-1)+1\cdot (z-4)=0 \Leftrightarrow 2x-4y+z+6=0.$$
			\itemch Ta có $\left[\vec{u_1};\vec{u_2}\right]= (-11;-7;-5)$, do đó mặt phẳng $(P)$ có một véc-tơ pháp tuyến là $\vec{n}=(-11;-7;-5)$. \\
			Phương trình mặt phẳng $(P)$ đi qua điểm $N(2;-1;5)$ là $$-11\cdot (x-2)-7\cdot (y+1) -5\cdot (z-5)=0 \Leftrightarrow -11x-7y-5z+40=0.$$
			\itemch Do $(P) \parallel (Q)$ nên mặt phẳng $(P)$ nhận véc-tơ pháp tuyến của mặt phẳng $(Q)$ làm véc-tơ pháp tuyến. Vậy mặt phẳng $(P)$ đi qua $I(4;0;-7)$ và $\vec{n}_{P}=(2;1;-1)$ có phương trình là $$2\cdot (x-4)+1\cdot (y-0)-1\cdot (z+7)=0 \Leftrightarrow 2x+y-z-15=0.$$
			\itemch Mặt phẳng $(P)$ vuông góc với đường thẳng $\Delta$ nên véc-tơ chỉ phương của $\Delta$ cũng là véc-tơ pháp tuyến của $(P)$. Vậy $(P)$ đi qua điểm $K(-4;9;2)$ và $\vec{n}_{P}=(2;1;5)$ nên có phương trình là $$2\cdot (x+4) +1\cdot (y-9)+5\cdot (z-2)=0 \Leftrightarrow 2x+y+5z-11=0.$$
		\end{itemchoice}
	}
\end{ex}
\begin{ex}%[2H5V3-3]
	Trong không gian $Oxyz$, cho ba điểm $A(1;0;-1)$, $B(0;1;2)$, $C(-1;-2;3)$. Mỗi mệnh đề sau là đúng hay sai?
	\choiceTF
	{Phương trình mặt phẳng $(ABC)$ là $5x-y+2z+3=0$}
	{\True Phương trình đường thẳng $AC$ là $\dfrac{x-1}{-1}=\dfrac{y}{1}=\dfrac{x+1}{3}$ }
	{Phương trình mặt cầu đường kính $AC$ là $x^2+(y-1)^2+(z+1)^2=12$}
	{\True Phương trình mặt cầu có tâm $A$ và đi qua $B$ là $(x-1)^2+y^2+(z+1)^2=11$}
	\loigiai{
		\begin{itemchoice}
			\itemch Mặt phẳng $(ABC)$ nhận $\overrightarrow{AB}=(-1;1;3)$, $\overrightarrow{AC}=(-2;-2;4)$ làm cặp véc-tơ chỉ phương nên có véc-tơ pháp tuyến là $\overrightarrow{n}=\left[\overrightarrow{AB},\overrightarrow{AC}\right]=(10;-2;4)$.\\
			Mặt phẳng $(ABC)$ đi qua $A(1;0;-1)$ và nhận $\overrightarrow{n}=(10;-2;4)$ làm một véc-tơ pháp tuyến nên có phương trình\\
			$10(x-1)-2y+4(z+1)=0\Leftrightarrow 10x-2y+4z-6=0\Leftrightarrow 5x-y+2z-3=0$.
			\itemch Đường thẳng $AB$ đi qua điểm $A(1;0;-1)$ nhận véc-tơ $\overrightarrow{AB}=(-1;1;3)$ làm một véc-tơ chỉ phương nên có phương trình là $\dfrac{x-1}{-1}=\dfrac{y}{1}=\dfrac{z+1}{3}$.
			\itemch Mặt cầu có tâm $I\left(0;-1;1\right)$ là trung điểm của đường kính $AC$.\\
			Bán kính $R=\dfrac{AC}{2}=\dfrac{\sqrt{(-2)^2+(-2)^2+4^2}}{2}=\sqrt{6}$.\\
			Phương trình mặt cầu là $x^2+(y+1)^2+(z-1)^2=6$.
			\itemch Mặt cầu có tâm $A(1;0;-1)$ và đi qua $B(0;1;2)$.\\
			Bán kính $R=AB=\sqrt{(-1)^2+1^2+3^2}=\sqrt{11}$.\\
			Phương trình mặt cầu là $(x-1)^2+y^2+(z+1)^2=11$.
		\end{itemchoice}
	}
\end{ex}
\begin{ex}%[2H5V3-3]
	Mỗi mệnh đề sau là đúng hay sai?
	\choiceTF
		{\True Mặt cầu $(S)$ có tâm $I(4;-2;1)$ và bán kính $R=9$ là $(x-4)^2+(y+2)^2+(z-1)^2=81$}
		{Mặt cầu $(S)$ có tâm $I(3;2;0)$ và đi qua điểm $M(2;4;-1)$ là $(S)$ là $(x+3)^2+(y+2)^2+z^2=6$}
		{Mặt cầu $(S)$ có đường kính là đoạn thẳng $AB$ với $A(1;2;0)$ và $B(-1;0;4)$ là $x^2+(y-1)^2+(z-2)^2=32$}
		{\True Mặt cầu $(S)$ mặt cầu tâm $I(-1 ; 2 ; 3)$ và tiếp xúc với mặt phẳng $(P)\colon 2 x-y-2 z+1=0$ có phương trình là $ (x+1)^2+(y-2)^2+(z-3)^2=9$}
	\loigiai{
		\begin{itemchoice}
			\itemch Phương trình mặt cầu $(S)$ có tâm $I(4;-2;1)$ và bán kính $R=9$ là $$(x-4)^2+(y+2)^2+(z-1)^2=81.$$
			\itemch Ta có $R=IM=\sqrt{(2-3)^2+(4-2)^2+(-1-0)^2} = \sqrt{6}$.\\
			Phương trình mặt cầu $(S)$ là $(x-3)^2+(y-2)^2+z^2=6$.
			\itemch Gọi $I$ là trung điểm $AB$, ta có $I(0;1;2)$, điểm $I$ chính là tâm của mặt cầu đường kính $AB$.\\
			Ta có $AB=\sqrt{(-1-1)^2+(0-2)^2+(4-0)^2}=2\sqrt{6}$. Suy ra $R= \dfrac{AB}{2}=\sqrt{6}$.\\
			Phương trình mặt cầu $(S)$ là $$(x-0)^2+(y-1)^2+(z-2)^2=6 \Leftrightarrow x^2+(y-1)^2+(z-2)^2=6.$$
			\itemch Vì  mặt cầu $(S)$ mặt cầu tâm $I(-1 ; 2 ; 3)$ và tiếp xúc với mặt phẳng $(P)$
			nên bán kính mặt cầu $$ R=\mathrm{d}(I,(P))=\dfrac{|2 \cdot (-1)-2-2 \cdot 3+1|}{\sqrt{2^2+(-1)^2+2^2}} =3.$$
			Phương trình mặt cầu tâm $ I $ tiếp xúc với $ (P) $ là
			$ (x+1)^2+(y-2)^2+(z-3)^2=9. $
		\end{itemchoice}
	}
\end{ex}
\begin{ex}%[2H5V2-7]
	Hình bên dưới minh họa đường bay của một chiếc trực thăng $H$ cất cánh từ một sân bay. Xét hệ trục tọa độ $Oxyz$ có gốc tọa độ $O$ là chân tháp điều khiển của sân bay; trục $Ox$ là hướng đông, trục $Oy$ là hướng Bắc và trục $Oz$ là trục thẳng đứng, đơn vị trên mỗi trục là kilômét.\\
	Trực thăng cất cánh từ điểm $G$. véc-tơ $\vec{r}$ chỉ vị trí của trực thăng tại thời điểm $t$ phút sau khi cất cánh $(t\ge 0)$ có tọa độ là $\vec{r}=(1+t; 0,5+2t; 2t)$.
	\begin{center}
		\begin{tikzpicture}
			\coordinate (O) at (0,0);
			\coordinate (z) at (0,3);
			\coordinate (x) at (2.5,-1.5);
			\coordinate (y) at (5,2);
			\coordinate (Oxy) at ($(y)+(2.5,-1.5)$);
			
			\path[fill=cyan!15] (O)--(x)--(Oxy)--(y)--cycle;
			\coordinate (G) at (1.5,0);
			\coordinate (F) at (4,0);
			\coordinate (F') at (5,0);
			\coordinate (H) at ($(F)+(0,3.5)$);
			\draw[->] (G)--(H);
			\draw[dashed] (G)--(F') (H)--(F);
			\draw ($(F)+(-0.2,0)$)|-($(F)+(0,0.2)$);
			\begin{scope}
				\clip (F)--(G)--(H);
				\draw (G) circle(0.4);
			\end{scope}
			\draw[->] (O)--(x);
			\draw[->] (O)--(y);
			\draw[->] (O)--(z);
			\foreach \x in{G,H,F}{
				\fill (\x) circle(1pt);
			}
			\path 
			node at (O) [below]{$O$}
			node at (z) [left]{$z$}
			node at (x) [below]{$x$ (Đ)}
			node at (y) [above right]{$y$ (B)}
			node at (G) [below]{$G$}
			node at (F) [below]{$F$}
			node at (H) [right]{$H$}
			node at ($(G)+(30:0.6)$) {$\theta$}
			;
		\end{tikzpicture}
	\end{center}
	Mỗi mệnh đề sau là đúng hay sai?
	\choiceTF
		{Góc $\theta$ mà đường bay tạo với phương ngang làm tròn đến phút là $45^\circ$}
		{\True Phương trình đường thẳng $GF$, trong đó $F$ là hình chiếu của điểm $H$ lên mặt phẳng $(Oxy)$ là $\heva{&x=1+t\\&y=0,5+2t\\&z=0}$, $t\in \mathbb{R}$}
		{\True Trực thăng bay vào mây ở độ cao $2$ km. Tọa độ điểm mà máy bay trực thăng bắt đầu đi vào đám mây là $(2;2{,}5;2)$}
		{Giả sử một đỉnh núi nằm ở điểm $M(5;4,5;3)$. Khi $HM$ vuông góc với đường bay $GH$ thì khoảng cách từ máy bay trực thăng đến đỉnh núi tại thời điểm đó là $3$km}
	\loigiai{
		\begin{itemchoice}
			\itemch Ta có phương trình tham số của đường thẳng $GH$ là $\heva{&x=1+t \\&y=0,5+2t\\ & z=2t}$, $t\in \mathbb{R}$.\\
			\allowdisplaybreaks
			\begin{eqnarray*}
			\sin \left(d; (Oxy)\right) &=& \left|\cos\left(\vec{u}_{d}; \vec{n}_{Oxy}\right)\right|\\
			&=& \dfrac{\left|\vec{u}_{d}\cdot \vec{n}_{Oxy}\right|}{\left|\vec{u}_{d}\right| \cdot \left| \vec{n}_{Oxy}\right|}\\
			&=& \dfrac{\left|1\cdot 0 + 2\cdot 0 + 2\cdot 1\right|}{\sqrt{1^2+2^2+2^2} \cdot \sqrt{0^2+0^2+1^2}}=\dfrac{2}{3}.	
			\end{eqnarray*}
			Suy ra $\left(d;(Oxy)\right)\approx 41^{\circ} 49'$. \\
			Vậy góc $\theta$ mà đường bay tạo với phương ngang xấp xỉ $41^{\circ} 49'$.
			\itemch Ta có điểm $M\left(2;2,5;4\right)$ nằm trên đường thẳng $GH$.\\ 
			Hình chiếu của $M$ lên mặt phẳng $(Oxy)$ là $N(2;2,5;0)$ và điểm $N$ nằm trên đường thẳng $GF$.\\
			Vậy $GF$  đi qua $G(1;0,5;0)$ và nhận véc-tơ $\vec{GN}=(1;2;0)$ làm véc-tơ chỉ phương, do đó có phương trình là $\heva{&x=1+t\\&y=0,5+2t\\&z=0}$, $t\in \mathbb{R}$.
			\itemch Trực thăng bay vào mây ở độ cao $2$ km, khi đó ta có 
			$2t=2 \Leftrightarrow t=1$. Tọa độ điểm mà máy bay bắt đầu đi vào mây là $(2;2{,}5;2)$.
			\itemch Gọi $H(1+t; 0,5+2t;2t)$ là vị trí cần tìm. Ta có $\vec{HM}=(4-t; 4-2t; 3-2t)$. \\
			Ta có $HM \perp GH$ suy ra $$\vec{HM}\cdot \vec{u}_{GH} = 0 \Leftrightarrow (4-t)\cdot 1 + (4-2t)\cdot 2 + (3-2t) \cdot 2 =0 \Leftrightarrow t=2.$$
			Vậy $H(3;4,5;4)$.\\
			Khoảng cách từ trực thăng đến đỉnh núi là $$MH = \sqrt{(3-5)^2+(4,5-4,5)^2+(4-3)^2} =\sqrt{5} \approx 2,236 \text{ km.}$$ 
			\itemch  
			\begin{center}
				\begin{tikzpicture}
					\coordinate (O) at (0,0);
					\coordinate (z) at (0,3);
					\coordinate (x) at (2.5,-1.5);
					\coordinate (y) at (5,2);
					\coordinate (Oxy) at ($(y)+(2.5,-1.5)$);
					
					\path[fill=cyan!15] (O)--(x)--(Oxy)--(y)--cycle;
					\coordinate (G) at (1.5,0);
					\coordinate (F) at (4,0);
					\coordinate (F') at (5,0);
					\coordinate (H) at ($(F)+(0,3.5)$);
					\draw[->] (G)--(H);
					\draw[dashed] (G)--(F') (H)--(F);
					
					\draw ($(F)+(-0.2,0)$)|-($(F)+(0,0.2)$);
					\begin{scope}
						\clip (F)--(G)--(H);
						\draw (G) circle(0.4);
					\end{scope}
					\draw[->] (O)--(x);
					\draw[->] (O)--(y);
					\draw[->] (O)--(z);
					\foreach \x in{G,H,F}{
						\fill (\x) circle(1pt);
					}
					\path 
					node at (O) [below]{$O$}
					node at (z) [left]{$z$}
					node at (x) [below]{$x$ (Đ)}
					node at (y) [above right]{$y$ (B)}
					node at (G) [below]{$G$}
					node at (F) [below]{$F$}
					node at (H) [right]{$H$}
					node at ($(G)+(30:0.6)$) {$\theta$}
					;
				\end{tikzpicture}
			\end{center}
		\end{itemchoice}
	}
\end{ex}


\Closesolutionfile{ans}
\indapan{3}{ans/ans-0-B15-DS}

\Opensolutionfile{ans}[ans/ans-0-B15-KQ]
\caukq
\begin{ex}%[2H2H2-2]
	Cho 3 điểm $A(2;-1;5)$, $B(5;-5;7)$, $M(x;y;1)$. Khi $A$, $B$, $M$ thẳng hàng thì $2x+y$ bằng?
	\shortans{$1$}
	\loigiai{
		Ta có $\overrightarrow{AB}=(3;-4;2)$, $\overrightarrow{AM}=(x-2;y+1;-4)$.\\
		$A$, $B$, $M$ thẳng hàng $\Leftrightarrow\overrightarrow{AB}$, $\overrightarrow{AM}$ cùng phương $\Leftrightarrow\dfrac{x-2}{3}=\dfrac{y+1}{-4}=\dfrac{-4}{2}\Leftrightarrow\heva{&x=-4\\&y=7.}$\\
		Vậy $2x+y=-8+7=1$
	}
\end{ex}
\begin{ex}%[2H5H1-3]
	Trong không gian $Oxyz$, cho điểm $A(1;2;-5)$. Gọi $M$, $N$, $P$ lần lượt là hình chiếu của điểm $A$ lên các trục $Ox$, $Oy$, $Oz$. Mặt phẳng $(MNP)$ có dạng phương trình $10x+by+cx+d=0$. Khi đó $b+c+d$ bằng
	\shortans{$-7$}
	\loigiai{Ta có $M(1;0;0)$, $N(0;2;0)$, $P(0;0-5)$. Mặt phẳng $(MNP)$ có phương trình là $$ x+\dfrac{y}{2}-\dfrac{z}{5}=1\Leftrightarrow 10x+5y-2z-10=0.$$
	Suy ra $b=5$, $c=-2$ và $d=-10$. Vậy $b+c+d=-7$.
	}
\end{ex}
\begin{ex}%[2H5V2-3]
	Trong không gian $Oxyz$, cho điểm $M(4;1;1)$ và mặt phẳng $(P) \colon 3x+y-z-1=0$. Hình chiếu vuông góc $H(a;b;c)$ của điểm $M$ trên mặt phẳng $(P)$. Khi đó $a-b+c$ bằng
	\shortans{$3$}
	\loigiai{	\immini
		{
			Gọi $ d $ là đường thẳng qua $M$ và vuông góc với $(P)$.\\
			Ta có $ \vv{u}_{d}=\vv{n}_{P}=(3;1;-1) $.\\
			Khi đó phương trình đường thẳng $ d\colon \heva{& x=4+3t \\ & y=1+t \\ & z=1-t.} $ \\
			Gọi $H = d \cap (P)$. Do $ H \in d \Rightarrow H(4+3t;1+t;1-t) \in (P)$.\\
		}
		{
			\begin{tikzpicture}[scale=1.2,font=\footnotesize,dot/.style={circle,inner sep=1pt,fill,label={#1},name=#1},
				extended line/.style={shorten >=-#1,shorten <=-#1},
				extended line/.default=1cm]
				\path 
				(0,0) coordinate (O)
				(2.5,0) coordinate (B)
				(B)--(O)--([turn]-130:1.1) coordinate (D)
				($(D)+(B)-( O)$) coordinate (C)
				($(O)!0.5!(C)$) coordinate (H)
				($(H)+(0,1)$) coordinate (M)node[shift={(0:10mm)}]{$M(4;1;1)$}
				($(H)+(0,1.5)$) coordinate (H1)
				($(M)!1.6!(H)$) coordinate (H2)
				($(M)!2.5!(H)$) coordinate (H3)
				
				($(H)!0.5!(O)$) coordinate (A)
				($(A)+(0,1)$) coordinate (A1)node[red,shift={(180:4.5mm)}]{$\vec{n}_{(P)}$}
				;
				\fill (M) circle (1pt);
				\draw (H1)--(H) (H3)--(H2);
				\draw[dashed](H)--(H3);
				\draw[->,>=stealth,red,thick](A)--(A1);
				\begin{scope}
					\clip (D)--(O)--(B);
					\draw[fill=cyan!60,opacity=0.7] (O) circle(0.6cm)node[black,shift={(25:4mm)}]{$P$};
				\end{scope}
				\draw (O)--(B)--(C)--(D)--cycle  ;
				\foreach \p/\r in {H/0}
				\fill (\p) circle (1pt) node[shift={(\r:3mm)}]{$\p$};
				
				%\foreach \p/\r in {H/90,M/90,H1/90,H2/90,H3/90,A1/90,A/90}
				%\fill (\p) circle (1pt) node[shift={(\r:3mm)}]{$\p$};
			\end{tikzpicture}
		}
		Thay tọa độ điểm $H$ vào phương trình mặt phẳng $(P)$, ta được
		\begin{equation*}
			3(4+3t)+1+t-(1-t)=0 \Leftrightarrow t=-1 \Rightarrow H(1;0;2) .
		\end{equation*}
		Vậy hình chiếu vuông góc của điểm $ M $ trên mặt phẳng $ (P) $ là điểm $ H(1;0;2)$.\\
		Vậy $a-b+c=1-0+2=3$
	}
\end{ex}
\begin{ex}%[2H5H1-3]
	Trong không gian $Oxyz$, cho mặt phẳng $(P)\colon 2x-y+2z-4=0$. Phương trình mặt phẳng $(Q)$ song song với $(P)$ và cách $(P)$ một khoảng bằng $1$ có dạng $2x-y+2z+m=0$. Khi đó $m$ bằng mấy?
	\shortans{$2$}
	\loigiai{Vì $(Q)$ song song với $(P)$ nên $(Q)$ có dạng $ 2x-y+2z+D=0$ với $ D\neq -4$.\\
		Lấy $M(0;-1;0)\in (P)$.\\
		Ta có 
		\allowdisplaybreaks
		\begin{eqnarray*}
			&& \mathrm{d}(P), (Q))=1\\
			&\Leftrightarrow&  \mathrm{d}(M, (Q))=1\\
			&\Leftrightarrow& \dfrac{|1+D|}{\sqrt{2^2+(-1)^2+2^2}}=1\\
			&\Leftrightarrow& |1+D|=3\\
			&\Leftrightarrow& \hoac{&D=2\\&D=-4 \text{ (loại)}.}
		\end{eqnarray*}
		Vậy phương trình mặt phẳng là $(Q)\colon 2x-y+2z+2=0$.\\
		Vậy $m=2$.}
\end{ex}


\begin{ex}
	Trong không gian $Oxyz$ cho điểm $ A(1;2;1) $ và hai đường thẳng $d_{1}\colon \dfrac{x-1}{1}=\dfrac{y-2}{-1}=\dfrac{z}{1}$; $d_{2}\colon \dfrac{x+1}{2}=\dfrac{y-1}{1}=\dfrac{z+3}{-2}$. Đường thẳng $ d $ đi qua $ A $, cắt $ d_1 $ và vuông góc với $ d_2 $ có véc-tơ chỉ phương $\overrightarrow{u}=(2;a;b)$. khi đó $a-b$ bằng . 
	\shortans{$-3$}
	\loigiai{
		Gọi $B(1+b;2-b;b)\in d_1 \Rightarrow \overrightarrow{AB}=(b;-b;b-1) $.\\ 
		Vì $ AB \perp d_2 $ nên $ \overrightarrow{AB}\cdot \overrightarrow{u}_{d_2}=0 \Leftrightarrow 2b-b-2(b-1)=0 \Leftrightarrow b=2$.\\
		Do đó $\overrightarrow{u}= \overrightarrow{AB}=(2;-2;1)$. Suy ra $a=-2$ và $b=1$.\\
		Vậy $a-b=-2-1=-3$. 
	}
\end{ex}
\begin{ex}%[2H5C3-4]
	Trong không gian với hệ tọa độ $Oxyz$, đài kiểm soát không lưu sân bay có tọa độ $O(0;0;0)$, mỗi đơn vị trên trục ứng với $1$ km. Máy bay bay trong phạm vi cách đài kiểm soát $417$ km sẽ hiển thị trên màn hình ra đa. Một máy bay đang ở vị trí $A\left(-688;-185;8\right)$, chuyển động theo đường thẳng $d$ có véc-tơ chỉ phương là $\vec{u}=(91;75;0)$ và hướng về đài kiểm soát không lưu.
	\begin{center}
		\begin{tikzpicture}
			\def\r{3}
			\def \x{3}  %bán kính trục lớn elip
			\def \y{1.2} 
			\coordinate (o1) at (-\r,0);
			\coordinate (o2) at (\r,0);
			\coordinate (O) at (0,0);
			
			% ve mp
			\coordinate (P) at (-5,2);
			\coordinate (Q) at ($(P)+(240:6)$);
			\coordinate (R) at ($(Q)+(12,0)$);
			\coordinate (S) at ($(R)+(60:6)$);
			\fill[cyan!20] (P)--(Q)--(R)--(S)--cycle;
			
			
			\draw (O) circle(\r); 
			\path[fill=orange!20] (O) circle(\r); 
			\fill[orange!80] (o2) arc (0:-180:\x cm and \y cm) arc (180:360:\x cm and \x cm);
			\draw[dashed] (o2) arc (0:180:\x cm and \y cm);
			\draw (o2) arc (0:-180:\x cm and \y cm);
			
			% truc oz
			\coordinate (Oz) at (0,4.5);
			\coordinate (Oz') at (0,2.5);
			\coordinate (Oz'') at (0,-3.5);
			\draw[dashed] (Oz'')--(Oz');
			\draw[->] (Oz')--(Oz);
			% tru  ox
			\def\t{240}
			\coordinate (Ox) at ({\x*cos(\t)},{\y*sin(\t)});
			\coordinate (Ox') at ($(O)!2!(Ox)$);
			\coordinate (Ox'') at ($(Ox)!2.2!(O)$);
			\draw[dashed] (O)--(Ox);
			\draw[->] (Ox)--(Ox');
			\draw[dashed] (O)--(Ox'');
			% truc oy
			\coordinate (O4) at (\r,0);
			\coordinate (O5) at (1.5*\r,0);
			\coordinate (O2) at (-\r,0);
			\coordinate (O1) at (-1.5*\r,0);
			\draw (O1)--(O2);
			\draw[dashed] (O2)--(O4);
			\draw[->] (O4)--(O5);
			% duong d
			\coordinate (A) at (-4.52,1.7);
			\coordinate (A3) at (4.5,0.3);
			\coordinate (A1) at ($(A)!0.25!(A3)$);
			\coordinate (A2) at ($(A)!0.82!(A3)$);
			\draw (A)--(A1);
			\draw[dashed] (A1)--(A2);
			\draw (A2)--(A3);
			\foreach \x in {O}{
				\fill (\x) circle(1pt);
			}
			\path 
			node at (O) [above left]{$O$}
			node at (A) [above] {$A$}
			node at (A1) [above]{B}
			node at (A2) [above right]{$C$}
			node at (A3) [above]{$d$}
			node at (Oz) [left]{$z$}
			node at (O5) [below]{$y$}
			node at (Ox')[below]{$x$}
			;
		\end{tikzpicture}
	\end{center}
	Tọa độ của vị trí sớm nhất mà máy bay xuất hiện trên màn hình ra đa là $M(a;b;c)$. Khi đó $a+b+c$ bằng
	\shortans{$-458$}
	\loigiai{
		Đường thẳng $d$ đi qua $A(-688;-185;8)$, có véc-tơ chỉ phương là $\vec{u}=(91;75;0)$ nên có phương trình tham số là $\heva{& x=-688 +91t \\ &y=-185+75t\\ &z=8}, t\in \mathbb{R}$.\\
			Mặt cầu $(S)$ tâm $O(0;0;0)$ bán kính $R=417$ có phương trình $x^2+y^2+z^2=173\,889$.\\
			Tọa độ giao điểm của đường thẳng $d$ và mặt cầu $(S)$ là nghiệm của hệ 
			\begin{eqnarray*}
				\heva{&x=-688 +91t \\ &y=-185+75t\\ &z=8 \\& x^2+y^2+z^2=173\,889} 
				&\Leftrightarrow & 
				\heva{&x=-688 +91t \\ &y=-185+75t\\ &z=8 \\& (-688+91t)^2+(-185+75t)^2+8^2=173\,889} \\
				& \Leftrightarrow & 	\heva{&x=-688 +91t \\ &y=-185+75t\\ &z=8 \\& \hoac{&t=8 \\&t=3.}}  
			\end{eqnarray*}
			Với $t=8 \Rightarrow $ tọa độ giao điểm thứ nhất là $M_1 (40;415;8)$.\\
			Với $t=3 \Rightarrow $ tọa độ giao điểm thứ hai là $M_2 (-506; 40;8)$.\\
			Ta có $AM_1 = \sqrt{889984}$ và $AM_2 = \sqrt{76869}$, suy ra $AM_2 < AM_1$. Vậy tọa độ của vị trí sớm nhất mà máy bay xuất hiện trên màn hình ra đa là $M_2=(-506;40;8)$.\\
			Điểm $M_2$ chính là điểm $B$ trên hình vẽ.\\
			Suy ra $a=-506$, $b=40$ và $c=8$. Vậy $a+b+c=-458$.
	}
\end{ex}
\Closesolutionfile{ans}
\indapan{6}{ans/ans-0-B15-KQ}